\section{Training LIMA}
\label{sec:training}

We train LIMA (Less Is More for Alignment) using the following protocol.
Starting from LLaMa 65B~\citep{touvron2023llama}, we fine-tune on our 1,000-example alignment training set.
To differentiate between each speaker (user and assistant), we introduce a special end-of-turn token (EOT) at the end of each utterance; this token plays the same role as EOS of halting generation, but avoids conflation with any other meaning that the pretrained model may have imbued into the preexisting EOS token.

We follow standard fine-tuning hyperparameters:
we fine-tune for 15 epochs using AdamW~\citep{loshchilov2017decoupled} with $\beta_1=0.9, \beta_2=0.95$, and weight decay of $0.1$.
Without warmup steps, we set the initial learning rate to $1e-5$ and linearly decaying to $1e-6$ by the end of training.
The batch size is set to 32 examples (64 for smaller models), and texts longer than 2048 tokens are trimmed.
One notable deviation from the norm is the use of residual dropout; we follow~\citet{ouyang2022training} and apply dropout over residual connections, starting at $p_d = 0.0$ at the bottom layer and linearly raising the rate to $p_d = 0.3$ at the last layer ($p_d = 0.2$ for smaller models).
We find that perplexity does not correlate with generation quality, and thus manually select checkpoints between the 5th and the 10th epochs using the held-out 50-example development set.\footnote{See Appendix~\ref{sec:ppl} for a more detailed study comparing validation perplexity and generation quality.}
